\section{Introduction}

% IEEE target for 7-page paper: ~1.0--1.2 pages.
% This section should define the problem, scope, and contribution clearly.
% Suggested flow: motivation, gap, contribution, and roadmap.

AI-assisted tools are increasingly used in cybersecurity workflows for triage, reporting,
explanation, and research support. This raises an operational question for defenders:
what measurable differences, if any, emerge when analysis work is performed under
documented AI-assisted conditions versus baseline conditions.

In this paper, we adopt a defensive, controlled research framing. The objective is not
to infer provenance of malware samples or describe offensive construction workflows. Instead,
we compare measurable outcomes under two known conditions and report differences that are
useful for detection engineering and analyst workflows.

The contributions are:
\begin{itemize}
  \item A defensive, conference-style scaffold for studying AI-assisted cybersecurity workflows.
  \item A clear condition-level methodology (baseline vs. AI-assisted) with reproducible metadata tracking.
  \item A structured set of static, dynamic, and detection-oriented metrics for comparative analysis.
\end{itemize}

The paper is organized as follows: Section II summarizes background and related work,
Section III states the research questions and study design, Section IV details the method,
Section V presents results, Section VI discusses implications and limitations, and
Section VII concludes with future research directions.
